\documentclass{ximera}
%herschreven 23 september 2026
\title{Elektrisch veld van twee bronladingen}
\author{Daan Lipkens}

\begin{document}
\begin{abstract}
    
\end{abstract}
\maketitle
\section{Dipoolveld}
Een dipool bestaat uit een positieve en een negatieve lading van exact dezelfde grootte op een kleine afstand van elkaar en maakt een dipoolveld, zoals te zien in figuur \ref{fig:veldlijnen_dipool}. Het aantal lijnen dat bij de positieve lading vertrekt, is gelijk aan het aantal lijnen dat op de negatieve lading toekomt. Tussen de ladingen in is de dichtheid van de lijnen hoog, wat wijst op een sterk veld.

%---------------------------
% AFBEELDING: Figuur 23.20 (Veldlijnen dipool)
\begin{figure}[ht!]
    \captionsetup{width=0.85\textwidth}
    \begin{center}
        \includegraphics[width=0.5\textwidth]{pic5/elektriciteit/veld/veldlijnen_dipool.png}
        % TIP: Gebruik hier Figure 23.20 (a)
    \end{center}
    \caption{Het elektrische veldlijnenpatroon van een dipool. De veldlijnen vertrekken bij de positieve lading en eindigen op de negatieve lading.}
    \label{fig:veldlijnen_dipool}
\end{figure}
%---------------------------


\section{Tussen twee gelijksoortige ladingen}
Bij twee identieke positieve ladingen vertrekken uit beide bronnen evenveel veldlijnen, zoals te zien op figuur \ref{fig:veldlijnen_gelijk}. Tussen de ladingen in 'botsen' de velden en buigen ze af naar boven en onder. Exact in het midden tussen de ladingen heffen de velden elkaar op en bevindt zich een punt waar het resulterende veld nul is.

%---------------------------
% AFBEELDING: Figuur 23.21 (Veldlijnen 2 positieve ladingen)
\begin{figure}[ht!]
    \captionsetup{width=0.85\textwidth}
    \begin{center}
        \includegraphics[width=0.5\textwidth]{pic5/elektriciteit/veld/veldlijnen_gelijk.png}
        % TIP: Gebruik hier Figure 23.21 (a)
    \end{center}
    \caption{Veldlijnen van twee gelijke positieve ladingen. In het exacte midden is het veld nul. (De plekken $A$, $B$ en $C$ worden gebruikt in de vraag \ref{v:veld2ladingen}.)}
    \label{fig:veldlijnen_gelijk}
\end{figure}
%---------------------------
%———————————————————————————————————————
% VRAAG: grootte van elektrisch veld op basis van veldlijnen
\begin{question}\label{v:veld2ladingen}
    Rangschik de grootte van het elektrisch veld in punten $A$, $B$ en $C$ van figuur \ref{fig:veldlijnen_gelijk} (grootste eerst).\\

    Rankschikking elektrische velden van groot naar klein:
    \wordChoice{\choice[correct]{A}\choice{B}\choice{C}}\qquad\wordChoice{\choice{A}\choice[correct]{B}\choice{C}}\qquad\wordChoice{\choice{A}\choice{B}\choice[correct]{C}}

    \begin{hint}
        Hoe dichter veldlijnen op elkaar staan, hoe groter het veld op die locatie is.
    \end{hint}

    \begin{solution}\nl
        De dichtheid van de veldlijnen is het grootste bij punt $A$, hier is dus het elektrisch veld het grootste. In punt $C$, wat net tussen de twee positieve ladingen ligt, ontbreken de veldlijnen. Hier is het totale elektrisch veld gelijk aan nul. Dus $E_A > E_B > E_C$.
    \end{solution}
\end{question}
 %———————————————————————————————————————
\end{document}